% ==================================================================
% id:       exp-004
% topic:    תוחלת ושונות
% subtopic: תוחלת של פונקציה של משתנה מקרי
% solution: full
% verified: no
% ==================================================================
עבור $X\sim\Bin(n,p)$, חשבו את $\E[X(X-1)]$ בשתי דרכים:
\begin{enumerate}
\item (הדרך הארוכה) בעזרת הכלל לתוחלת של פונקציה של משתנה מקרי (כלל \textenglish{LOTUS}). מתחילים כך:
$\E[X(X-1)]=\sum_{k=0}^nk(k-1)\PP(X=k)$. אפשר להיעזר בזהות $k\binom nk=n\binom{n-1}{k-1}$.
\begin{solbox}
\unverified
לפי הכלל לתוחלת של פונקציה של משתנה מקרי:
\[
\begin{split}
\E[X(X-1)]&=\sum_{k=0}^nk(k-1)\PP(X=k)\overset{(1)}=\sum_{k=2}^nk(k-1)\binom nkp^k(1-p)^{n-k}\\
&\overset{(2)}=\sum_{k=2}^n(k-1)\,n\binom{n-1}{k-1}p^k(1-p)^{n-k}\\
&\overset{(3)}=n(n-1)\sum_{k=2}^n\binom{n-2}{k-2}p^k(1-p)^{n-k}
\\&=n(n-1)p^2\sum_{k=2}^n\binom{n-2}{k-2}p^{k-2}(1-p)^{(n-2)-(k-2)}\\
&\overset{(4)}=n(n-1)p^2\sum_{\ell=0}^{n-2}\binom{n-2}\ell p^\ell(1-p)^{n-2-\ell}\overset{(5)}=\ans{n(n-1)p^2}.
\end{split}
\]
הסברים למעברים:
(1) עבור $k=0,1$ המחוברים מתאפסים.
(2) הזהות $k\binom nk=n\binom{n-1}{k-1}$.
(3) אותה זהות שוב: $(k-1)\binom{n-1}{k-1}=(n-1)\binom{n-2}{k-2}$.
(4) החלפת אינדקס $\ell=k-2$.
(5) זה סכום כל ההסתברויות של משתנה $\Bin(n-2,p)$, ולכן הוא שווה $1$.
\end{solbox}
\item (הדרך הקצרה) בעזרת לינאריות התוחלת והתוחלת והשונות של $X$.
\begin{solbox}
\unverified
לפי לינאריות התוחלת ו-$\E[X^2]=\Var[X]+(\E[X])^2$:
\begin{multline*}
\E[X(X-1)]=\E[X^2]-\E[X]=\Var[X]+(\E[X])^2-\E[X]\\
=np(1-p)+(np)^2-np=\ans{n(n-1)p^2}.
\end{multline*}
\end{solbox}
\end{enumerate}
